% !TEX program = xelatex
\documentclass[11pt,a4paper]{article}

% --- 中文支持 ---
\usepackage[UTF8,heading=true]{ctex}
\setCJKmainfont{SimSun}[BoldFont=SimHei,ItalicFont=KaiTi]
\setCJKsansfont{SimHei}
\setCJKmonofont{FangSong}

% --- 页面布局：窄边距，最大化画布 ---
\usepackage[top=12mm,bottom=12mm,left=12mm,right=12mm]{geometry}
\pagestyle{empty}

% --- 颜色系统 ---
\usepackage[dvipsnames,svgnames,x11names]{xcolor}
\definecolor{primary}{HTML}{1565C0}
\definecolor{primaryLight}{HTML}{42A5F5}
\definecolor{primaryDark}{HTML}{0D47A1}
\definecolor{accent}{HTML}{FF6F00}
\definecolor{accentLight}{HTML}{FFA726}
\definecolor{success}{HTML}{2E7D32}
\definecolor{successLight}{HTML}{66BB6A}
\definecolor{danger}{HTML}{C62828}
\definecolor{dangerLight}{HTML}{EF5350}
\definecolor{purple}{HTML}{6A1B9A}
\definecolor{purpleLight}{HTML}{AB47BC}
\definecolor{teal}{HTML}{00695C}
\definecolor{tealLight}{HTML}{26A69A}
\definecolor{bgPage}{HTML}{F5F7FA}
\definecolor{bgCard}{HTML}{FFFFFF}
\definecolor{gray80}{HTML}{424242}
\definecolor{gray60}{HTML}{757575}
\definecolor{gray40}{HTML}{BDBDBD}
\definecolor{gray20}{HTML}{E0E0E0}
\definecolor{ocean}{HTML}{0277BD}
\definecolor{oceanLight}{HTML}{4FC3F7}
\definecolor{amber}{HTML}{F57F17}

% --- TikZ ---
\usepackage{tikz}
\usetikzlibrary{arrows.meta,positioning,shadows.blur,calc,shapes.geometric,
  fit,backgrounds,decorations.pathreplacing,patterns,fadings,trees}

% --- 其他包 ---
\usepackage{fontawesome5}
\usepackage{graphicx}
\usepackage{setspace}
\usepackage{anyfontsize}
\usepackage{hyperref}
\hypersetup{colorlinks=false}
\usepackage{amsmath,amssymb}

% ==========================================
% 全局公用样式定义
% ==========================================
\tikzset{
  card/.style={
    rectangle, rounded corners=6pt, draw=gray20, line width=0.4pt, fill=bgCard,
    inner sep=8pt, align=center,
    blur shadow={shadow blur steps=6,shadow xshift=0.6pt,shadow yshift=-0.8pt,
      shadow blur radius=2pt,shadow opacity=12}
  },
  card/.default=50mm,
  iconCircle/.style={
    circle, minimum size=14mm, fill=#1, text=white,
    font=\Large, inner sep=0pt
  },
  iconCircle/.default=primary,
  bigIconCircle/.style={
    circle, minimum size=18mm, fill=#1, text=white,
    font=\LARGE, inner sep=0pt
  },
  bigIconCircle/.default=primary,
  phase/.style={
    rectangle, rounded corners=4pt, minimum width=145mm, minimum height=16mm,
    fill=#1, text=white, font=\sffamily\bfseries\normalsize,
    inner xsep=10pt, align=left,
    blur shadow={shadow blur steps=4,shadow xshift=0.4pt,shadow yshift=-0.6pt,
      shadow blur radius=1.5pt,shadow opacity=10}
  },
  phase/.default=primary,
  badge/.style={
    rectangle, rounded corners=2pt, minimum height=6mm,
    fill=#1, text=white, font=\sffamily\bfseries\tiny,
    inner xsep=5pt, inner ysep=1pt
  },
  badge/.default=primary,
  formulabox/.style={
    rectangle, rounded corners=8pt, draw=#1, line width=1.2pt,
    fill=#1!5, inner sep=10pt, align=center
  },
  formulabox/.default=primary,
}

\newcommand{\pagemark}[1]{%
  \begin{tikzpicture}[remember picture,overlay]
    \node[anchor=south east,font=\sffamily\tiny\color{gray40}]
      at ([xshift=-8mm,yshift=8mm]current page.south east) {#1};
  \end{tikzpicture}%
}

\begin{document}

% ============================================================
%  PAGE 1: 封面
% ============================================================
\begin{tikzpicture}[remember picture, overlay]
  % 渐变背景
  \shade[left color=primaryDark,right color=ocean]
    (current page.south west) rectangle (current page.north east);
  % 装饰圆
  \fill[white,opacity=0.04] ([xshift=5cm,yshift=-3cm]current page.north east) circle(10cm);
  \fill[white,opacity=0.03] ([xshift=-6cm,yshift=2cm]current page.south west) circle(8cm);
  \fill[white,opacity=0.02] ([xshift=2cm,yshift=5cm]current page.south west) circle(6cm);

  % 顶部标签
  \node[rectangle, rounded corners=3pt, fill=white!15, inner xsep=12pt, inner ysep=4pt,
    font=\sffamily\small\color{white!90}]
    at ([yshift=6cm]current page.center)
    {\faBook\quad 流体力学 $\cdot$ 期中复习系列};

  % 主标题
  \node[text width=16cm,align=center] at ([yshift=2.5cm]current page.center) {
    {\color{white}\sffamily\fontsize{42}{50}\selectfont\bfseries 第四章}\\[8pt]
    {\color{white}\sffamily\fontsize{28}{36}\selectfont 雷诺输运公式与连续性方程}
  };

  % 分割线
  \draw[white,opacity=0.4,line width=0.8pt]
    ([xshift=-6cm,yshift=-0.3cm]current page.center)--([xshift=6cm,yshift=-0.3cm]current page.center);

  % 副标题
  \node[text width=14cm,align=center] at ([yshift=-1.5cm]current page.center) {
    {\color{white!85}\sffamily\large 从系统到控制体的核心桥梁}\\[4pt]
    {\color{white!70}\sffamily\small 考试分值预估：8{\raise.17ex\hbox{$\scriptstyle\sim$}}12分\quad|\quad 填空 + 判断 + 计算题}
  };

  % 底部三个核心关键词卡片
  \node[rectangle,rounded corners=8pt,fill=white!12,
    minimum width=42mm,minimum height=20mm,text=white,align=center] (k1)
    at ([xshift=-5.2cm,yshift=-5cm]current page.center) {
    {\sffamily\bfseries\large \faRandom\; 系统 vs 控制体}\\[-1pt]
    {\sffamily\footnotesize 两种分析视角}
  };
  \node[rectangle,rounded corners=8pt,fill=white!12,
    minimum width=42mm,minimum height=20mm,text=white,align=center] (k2)
    at ([yshift=-5cm]current page.center) {
    {\sffamily\bfseries\large \faSuperscript\; 雷诺输运方程}\\[-1pt]
    {\sffamily\footnotesize 翻译工具}
  };
  \node[rectangle,rounded corners=8pt,fill=white!12,
    minimum width=42mm,minimum height=20mm,text=white,align=center] (k3)
    at ([xshift=5.2cm,yshift=-5cm]current page.center) {
    {\sffamily\bfseries\large \faWater\; 连续性方程}\\[-1pt]
    {\sffamily\footnotesize 质量守恒}
  };

  % 底部装饰
  \node[font=\sffamily\tiny\color{white!40}]
    at ([yshift=1.2cm]current page.south) {复习指南 v1.0};
\end{tikzpicture}
\pagemark{1/8}

\clearpage
% ============================================================
%  PAGE 2: 系统 vs 控制体 概念对比
% ============================================================
\begin{tikzpicture}[remember picture, overlay]
  % 浅灰背景
  \fill[bgPage] (current page.south west) rectangle (current page.north east);

  % 页眉
  \fill[primary] (current page.north west) rectangle ([yshift=-2cm]current page.north east);
  \node[anchor=west,font=\sffamily\LARGE\bfseries\color{white}]
    at ([xshift=2cm,yshift=-1cm]current page.north west) {\faRandom\; 核心概念：系统 vs 控制体};

  % ---- 左卡片：系统 ----
  \node[card, minimum width=82mm, minimum height=105mm] (sysCard)
    at ([xshift=-4.6cm,yshift=1.5cm]current page.center) {};

  \node[iconCircle=primary] at ([yshift=35mm]sysCard.center) {\faUsers};
  \node[font=\sffamily\bfseries\Large\color{primary}] at ([yshift=18mm]sysCard.center) {系统（体系）};
  \node[font=\sffamily\normalsize\color{gray80},text width=70mm,align=center]
    at ([yshift=5mm]sysCard.center) {始终由\textbf{同一批流体质点}组成的集合};

  \node[rectangle, rounded corners=4pt, fill=primary!8, draw=primary!30,
    text width=68mm, align=left, inner sep=6pt, font=\sffamily\small\color{gray80}]
    at ([yshift=-18mm]sysCard.center) {%
    \faCheck\; 质量始终不变\\[2pt]
    \faCheck\; 边界随流体运动变形\\[2pt]
    \faCheck\; 拉格朗日视角（盯人防守）\\[2pt]
    \faCheck\; 物理定律的天然对象
  };

  % ---- 右卡片：控制体 ----
  \node[card, minimum width=82mm, minimum height=105mm] (cvCard)
    at ([xshift=4.6cm,yshift=1.5cm]current page.center) {};

  \node[iconCircle=accent] at ([yshift=35mm]cvCard.center) {\faCube};
  \node[font=\sffamily\bfseries\Large\color{accent}] at ([yshift=18mm]cvCard.center) {控制体（CV）};
  \node[font=\sffamily\normalsize\color{gray80},text width=70mm,align=center]
    at ([yshift=5mm]cvCard.center) {流场中\textbf{某个确定的空间区域}};

  \node[rectangle, rounded corners=4pt, fill=accent!8, draw=accent!30,
    text width=68mm, align=left, inner sep=6pt, font=\sffamily\small\color{gray80}]
    at ([yshift=-18mm]cvCard.center) {%
    \faCheck\; 空间位置/形状通常固定\\[2pt]
    \faCheck\; 流体可穿越边界进出\\[2pt]
    \faCheck\; 欧拉视角（守株待兔）\\[2pt]
    \faCheck\; 工程分析的常用视角
  };

  % 中间连接箭头
  \node[font=\sffamily\bfseries\huge\color{gray40}] at ([yshift=1.5cm]current page.center) {\faArrowsAltH};

  % ---- 底部：真题高频考点 ----
  \node[rectangle, rounded corners=6pt, fill=danger!8, draw=danger!40,
    line width=0.8pt, minimum width=170mm, minimum height=35mm,
    align=left, inner sep=10pt] (examBox)
    at ([yshift=-5.8cm]current page.center) {};

  \node[anchor=north west, font=\sffamily\bfseries\normalsize\color{danger}]
    at ([xshift=5pt,yshift=-3pt]examBox.north west)
    {\faExclamationTriangle\; 真题高频考点};

  \node[anchor=north west, font=\sffamily\small\color{gray80}, text width=155mm]
    at ([xshift=5pt,yshift=-18pt]examBox.north west) {%
    \textbf{填空}："控制体定义为什么？" $\longrightarrow$ \textbf{流场中某个确定的空间区域}\\[3pt]
    \textbf{判断}："系统的表面通常会不断的变形" $\longrightarrow$ \textbf{对}（系统跟踪同一批质点，流体运动时表面自然变形）
  };
\end{tikzpicture}
\pagemark{2/8}

\clearpage
% ============================================================
%  PAGE 3: 雷诺输运方程（RTT）
% ============================================================
\begin{tikzpicture}[remember picture, overlay]
  \fill[bgPage] (current page.south west) rectangle (current page.north east);

  % 页眉
  \fill[ocean] (current page.north west) rectangle ([yshift=-2cm]current page.north east);
  \node[anchor=west,font=\sffamily\LARGE\bfseries\color{white}]
    at ([xshift=2cm,yshift=-1cm]current page.north west)
    {\faSuperscript\; 雷诺输运方程 (Reynolds Transport Theorem)};

  % 核心思想
  \node[rectangle, rounded corners=6pt, fill=ocean!10, draw=ocean!40,
    line width=0.8pt, minimum width=170mm, minimum height=18mm,
    font=\sffamily\normalsize\color{gray80}, text width=160mm, align=center, inner sep=6pt]
    at ([yshift=5.5cm]current page.center) {%
    力学定律是针对\textbf{系统}建立的，但我们用\textbf{控制体}分析流动\\[2pt]
    \textbf{雷诺输运方程}就是那个``\textbf{翻译工具}''——把系统视角 $\rightarrow$ 控制体视角
  };

  % 中心公式框
  \node[formulabox=primaryDark, minimum width=170mm, minimum height=32mm] (rttBox)
    at ([yshift=1.2cm]current page.center) {%
    {\color{primaryDark}\sffamily\bfseries\large 雷诺输运方程}\\[6pt]
    $\displaystyle \frac{DB_{\text{sys}}}{Dt} = \frac{\partial}{\partial t}\int_{CV} \rho\, b \; dV + \oint_{CS} \rho\, b \;(\vec{v} \cdot \vec{n}) \; dA$
  };

  % 三项含义 —— 三个卡片
  \node[card, minimum width=52mm, minimum height=50mm] (t1)
    at ([xshift=-5.5cm,yshift=-4cm]current page.center) {};
  \node[badge=primaryDark] at ([yshift=16mm]t1.center) {左端};
  \node[font=\sffamily\small\color{gray80}, text width=44mm, align=center]
    at ([yshift=2mm]t1.center) {$\displaystyle\frac{DB_{\text{sys}}}{Dt}$};
  \node[font=\sffamily\footnotesize\color{gray60}, text width=44mm, align=center]
    at ([yshift=-12mm]t1.center) {\textbf{系统}中物理量$B$\\的随体变化率};

  \node[card, minimum width=52mm, minimum height=50mm] (t2)
    at ([yshift=-4cm]current page.center) {};
  \node[badge=primary] at ([yshift=16mm]t2.center) {右端第一项};
  \node[font=\sffamily\small\color{gray80}, text width=44mm, align=center]
    at ([yshift=2mm]t2.center) {$\displaystyle\frac{\partial}{\partial t}\!\int_{CV}\!\rho b\,dV$};
  \node[font=\sffamily\footnotesize\color{gray60}, text width=44mm, align=center]
    at ([yshift=-12mm]t2.center) {控制体内$B$的\\当地变化率（储存量）};

  \node[card, minimum width=52mm, minimum height=50mm] (t3)
    at ([xshift=5.5cm,yshift=-4cm]current page.center) {};
  \node[badge=accent] at ([yshift=16mm]t3.center) {右端第二项};
  \node[font=\sffamily\small\color{gray80}, text width=44mm, align=center]
    at ([yshift=2mm]t3.center) {$\displaystyle\oint_{CS}\!\rho b(\vec{v}\!\cdot\!\vec{n})\,dA$};
  \node[font=\sffamily\footnotesize\color{gray60}, text width=44mm, align=center]
    at ([yshift=-12mm]t3.center) {$B$通过控制面的\\净通量（流出$-$流入）};

  % 箭头连接
  \draw[-{Stealth[length=5pt]},ocean,line width=1.2pt]
    ([yshift=-18mm]rttBox.south) -- ([yshift=27mm]t2.north);

  % 底部必背文字表述
  \node[rectangle, rounded corners=6pt, fill=amber!10, draw=amber!50,
    line width=0.8pt, minimum width=170mm, minimum height=20mm,
    inner sep=8pt, align=center] (memBox)
    at ([yshift=-8.5cm]current page.center) {};
  \node[anchor=north, font=\sffamily\small\color{gray80}, text width=155mm, align=center]
    at ([yshift=-3pt]memBox.north) {%
    {\color{amber}\bfseries\faLightbulb\; 真题必背文字表述}\\[4pt]
    系统中物理量的随体导数 = 控制体内该物理量对时间的偏导数 + 通过控制面的该物理量的净通量
  };

\end{tikzpicture}
\pagemark{3/8}

\clearpage
% ============================================================
%  PAGE 4: 连续性方程
% ============================================================
\begin{tikzpicture}[remember picture, overlay]
  \fill[bgPage] (current page.south west) rectangle (current page.north east);

  % 页眉
  \fill[success] (current page.north west) rectangle ([yshift=-2cm]current page.north east);
  \node[anchor=west,font=\sffamily\LARGE\bfseries\color{white}]
    at ([xshift=2cm,yshift=-1cm]current page.north west)
    {\faWater\; 连续性方程（质量守恒）};

  % 推导过程
  \node[rectangle, rounded corners=6pt, fill=success!8, draw=success!30,
    line width=0.8pt, minimum width=170mm, minimum height=22mm,
    inner sep=8pt, align=left, font=\sffamily\small\color{gray80}]
    at ([yshift=5.2cm]current page.center) {%
    {\color{success}\bfseries 推导思路：}令 $B=m$（质量），$b=B/m=1$ 代入RTT，由 $\frac{Dm_{\text{sys}}}{Dt}=0$ 得
  };

  % 积分形式
  \node[formulabox=success, minimum width=170mm, minimum height=22mm] (intBox)
    at ([yshift=2cm]current page.center) {%
    {\color{success}\sffamily\bfseries 积分形式}\qquad
    $\displaystyle \frac{\partial}{\partial t}\int_{CV} \rho \; dV + \oint_{CS} \rho (\vec{v} \cdot \vec{n}) \; dA = 0$
  };

  % 微分形式
  \node[formulabox=primaryDark, minimum width=170mm, minimum height=22mm] (diffBox)
    at ([yshift=-1.2cm]current page.center) {%
    {\color{primaryDark}\sffamily\bfseries 微分形式}\qquad
    $\displaystyle \frac{\partial \rho}{\partial t} + \nabla \cdot (\rho \vec{v}) = 0$
    \qquad{\color{gray60}\small\sffamily 即：密度对时间的偏导 + 密度与速度乘积的散度 = 0}
  };

  % 简化形式表格 —— 三个条件卡片
  \node[font=\sffamily\bfseries\large\color{gray80}]
    at ([xshift=-6cm,yshift=-3.8cm]current page.center) {\faFilter\; 各种简化形式};

  % 不可压缩
  \node[card, minimum width=52mm, minimum height=48mm] (s1)
    at ([xshift=-5.5cm,yshift=-6.5cm]current page.center) {};
  \node[badge=primary] at ([yshift=14mm]s1.center) {不可压缩流体};
  \node[font=\sffamily\normalsize\color{gray80}, text width=44mm, align=center]
    at ([yshift=0mm]s1.center) {$\nabla \cdot \vec{v} = 0$};
  \node[font=\sffamily\footnotesize\color{gray60}, text width=44mm, align=center]
    at ([yshift=-12mm]s1.center) {速度场散度为零\\与是否定常\textbf{无关}};

  % 定常可压缩
  \node[card, minimum width=52mm, minimum height=48mm] (s2)
    at ([yshift=-6.5cm]current page.center) {};
  \node[badge=accent] at ([yshift=14mm]s2.center) {定常可压缩流};
  \node[font=\sffamily\normalsize\color{gray80}, text width=44mm, align=center]
    at ([yshift=0mm]s2.center) {$\nabla \cdot (\rho\vec{v}) = 0$};
  \node[font=\sffamily\footnotesize\color{gray60}, text width=44mm, align=center]
    at ([yshift=-12mm]s2.center) {时间项消失\\但密度不能提出来};

  % 一维定常不可压
  \node[card, minimum width=52mm, minimum height=48mm] (s3)
    at ([xshift=5.5cm,yshift=-6.5cm]current page.center) {};
  \node[badge=teal] at ([yshift=14mm]s3.center) {一维定常不可压};
  \node[font=\sffamily\normalsize\color{gray80}, text width=44mm, align=center]
    at ([yshift=0mm]s3.center) {$A_1 V_1 = A_2 V_2$};
  \node[font=\sffamily\footnotesize\color{gray60}, text width=44mm, align=center]
    at ([yshift=-12mm]s3.center) {管道流的基石\\管细则速快};

  % 经典判断题陷阱
  \node[rectangle, rounded corners=6pt, fill=danger!8, draw=danger!40,
    line width=0.8pt, minimum width=170mm, minimum height=18mm,
    inner sep=6pt, align=center] (trapBox)
    at ([yshift=-10cm]current page.center) {};
  \node[font=\sffamily\small\color{gray80}, text width=160mm, align=center]
    at (trapBox.center) {%
    {\color{danger}\bfseries \faExclamationTriangle\; 经典陷阱}\;
    ``理想流体定常流动的连续性方程为 $\nabla\cdot\vec{v}=0$'' $\longrightarrow$
    {\color{danger}\bfseries 错！}$\nabla\cdot\vec{v}=0$ 是\textbf{不可压缩}流体的方程
  };
\end{tikzpicture}
\pagemark{4/8}

\clearpage
% ============================================================
%  PAGE 5: 工程应用——三大题型
% ============================================================
\begin{tikzpicture}[remember picture, overlay]
  \fill[bgPage] (current page.south west) rectangle (current page.north east);

  % 页眉
  \fill[accent] (current page.north west) rectangle ([yshift=-2cm]current page.north east);
  \node[anchor=west,font=\sffamily\LARGE\bfseries\color{white}]
    at ([xshift=2cm,yshift=-1cm]current page.north west)
    {\faTools\; 连续性方程的三大应用题型};

  % ---- 题型一 ----
  \node[card, minimum width=170mm, minimum height=52mm] (ty1)
    at ([yshift=4.2cm]current page.center) {};
  \node[bigIconCircle=primary] at ([xshift=-6.5cm,yshift=0mm]ty1.center) {\faCheckDouble};
  \node[anchor=west, font=\sffamily\bfseries\large\color{primary}]
    at ([xshift=-4.5cm,yshift=12mm]ty1.center) {题型一：验证流动是否连续};
  \node[anchor=west, font=\sffamily\small\color{gray80}, text width=120mm]
    at ([xshift=-4.5cm,yshift=-5mm]ty1.center) {%
    \textbf{方法}：对不可压缩流体，验证 $\frac{\partial u}{\partial x} + \frac{\partial v}{\partial y} + \frac{\partial w}{\partial z} = 0$\\[3pt]
    只需分别求偏导再相加，结果为0即满足连续性方程
  };

  % ---- 题型二 ----
  \node[card, minimum width=170mm, minimum height=52mm] (ty2)
    at ([yshift=-1.2cm]current page.center) {};
  \node[bigIconCircle=accent] at ([xshift=-6.5cm,yshift=0mm]ty2.center) {\faRocket};
  \node[anchor=west, font=\sffamily\bfseries\large\color{accent}]
    at ([xshift=-4.5cm,yshift=12mm]ty2.center) {题型二：管道变截面求流速};
  \node[anchor=west, font=\sffamily\small\color{gray80}, text width=120mm]
    at ([xshift=-4.5cm,yshift=-5mm]ty2.center) {%
    \textbf{核心公式}：$A_1 V_1 = A_2 V_2$，圆管面积 $A = \pi d^2/4$\\[3pt]
    直径缩小一半 $\rightarrow$ 面积缩至$1/4$ $\rightarrow$ 流速增大4倍
  };

  % ---- 题型三 ----
  \node[card, minimum width=170mm, minimum height=62mm] (ty3)
    at ([yshift=-6.2cm]current page.center) {};
  \node[bigIconCircle=teal] at ([xshift=-6.5cm,yshift=5mm]ty3.center) {\faWater};
  \node[anchor=west, font=\sffamily\bfseries\large\color{teal}]
    at ([xshift=-4.5cm,yshift=16mm]ty3.center) {题型三：水箱液面变化率（真题常考）};
  \node[anchor=west, font=\sffamily\small\color{gray80}, text width=120mm]
    at ([xshift=-4.5cm,yshift=-5mm]ty3.center) {%
    \textbf{标准三步法}：\\[2pt]
    \textbf{Step 1}：取整个容器为控制体\\[2pt]
    \textbf{Step 2}：储存项 = $\rho A_0 \frac{dh}{dt}$，通量项 = $\rho(A_2 V_2 - A_1 V_1)$\\[2pt]
    \textbf{Step 3}：合并得
    $\displaystyle\frac{dh}{dt} = \frac{A_1 V_1 - A_2 V_2}{A_0}$
  };
\end{tikzpicture}
\pagemark{5/8}

\clearpage
% ============================================================
%  PAGE 6: 例题——水箱液面变化率（含示意图）
% ============================================================
\begin{tikzpicture}[remember picture, overlay]
  \fill[bgPage] (current page.south west) rectangle (current page.north east);

  % 页眉
  \fill[teal] (current page.north west) rectangle ([yshift=-2cm]current page.north east);
  \node[anchor=west,font=\sffamily\LARGE\bfseries\color{white}]
    at ([xshift=2cm,yshift=-1cm]current page.north west)
    {\faPen\; 例题演示：水箱液面变化率};

  % 左侧示意图
  \begin{scope}[shift={([xshift=-4.6cm,yshift=2.5cm]current page.center)}]
    % 水箱外框
    \draw[gray80, line width=1.5pt] (-2.5,-3) rectangle (2.5,3);
    % 交叉线表示壁面
    \foreach \y in {-2.8,-2.4,...,2.8} {
      \draw[gray40, line width=0.3pt] (-2.7,\y) -- (-2.5,\y+0.2);
      \draw[gray40, line width=0.3pt] (2.5,\y) -- (2.7,\y+0.2);
    }
    % 水面
    \draw[primaryLight, line width=2pt] (-2.5,0) -- (2.5,0);
    \node[font=\tiny\color{primaryLight}] at (0,0.25) {$\sim\sim\sim\sim\sim\sim\sim\sim\sim\sim\sim\sim$};
    % 水区域
    \fill[primaryLight!15] (-2.5,-3) rectangle (2.5,0);
    % 空气区域
    \node[font=\sffamily\small\color{gray60}] at (0,1.8) {空气 $\rho_a$};
    \node[font=\sffamily\small\color{primary}] at (0,-1.5) {水 $\rho$};
    % 标注 H
    \draw[{Stealth}-{Stealth},gray80,line width=0.8pt] (3,-3) -- (3,3);
    \node[font=\sffamily\small\color{gray80},anchor=west] at (3.2,0) {$H$};
    % 标注 h
    \draw[{Stealth}-{Stealth},primary,line width=0.8pt] (3.8,-3) -- (3.8,0);
    \node[font=\sffamily\small\color{primary},anchor=west] at (4,{-1.5}) {$h$};
    % 标注 A0
    \draw[{Stealth}-{Stealth},gray60,line width=0.6pt] (-2.5,3.3) -- (2.5,3.3);
    \node[font=\sffamily\small\color{gray60}] at (0,3.7) {$A_0$};
    % 进水口
    \draw[success, line width=1.5pt, -{Stealth[length=5pt]}] (-4.5,-2) -- (-2.5,-2);
    \node[font=\sffamily\footnotesize\color{success},anchor=east] at (-4.6,-2) {$A_1,\,V_1$};
    \node[font=\sffamily\tiny\color{success},anchor=east] at (-4.6,-2.5) {进水};
    % 出水口
    \draw[danger, line width=1.5pt, -{Stealth[length=5pt]}] (2.5,-2) -- (4.5,-2);
    \node[font=\sffamily\footnotesize\color{danger},anchor=west] at (4.6,-2) {$A_2,\,V_2$};
    \node[font=\sffamily\tiny\color{danger},anchor=west] at (4.6,-2.5) {出水};
  \end{scope}

  % 右侧解题过程
  \node[anchor=north west, font=\sffamily\small\color{gray80}, text width=72mm, align=left]
    at ([xshift=1.5cm,yshift=5.5cm]current page.center) {%
    {\color{teal}\bfseries Step 1：选取控制体}\\[2pt]
    取整个水箱为控制体\\[6pt]
    {\color{teal}\bfseries Step 2：分析各项}\\[2pt]
    储存项（水+空气）：\\
    $(\rho-\rho_a)A_0\frac{dh}{dt}$\\[4pt]
    净流出通量：\\
    $\rho A_2 V_2 - \rho A_1 V_1$\\[6pt]
    {\color{teal}\bfseries Step 3：连续性方程}\\[2pt]
    $(\rho-\rho_a)A_0\frac{dh}{dt}$\\$+\;\rho A_2 V_2 - \rho A_1 V_1 = 0$
  };

  % 最终结果框
  \node[formulabox=teal, minimum width=76mm, minimum height=26mm]
    at ([xshift=4.2cm,yshift=-4.2cm]current page.center) {%
    {\color{teal}\sffamily\bfseries 最终结果}\\[4pt]
    $\displaystyle\frac{dh}{dt} = \frac{\rho(A_1 V_1 - A_2 V_2)}{(\rho - \rho_a)A_0}$
  };

  % 底部注释
  \node[rectangle, rounded corners=4pt, fill=successLight!12, draw=successLight!40,
    minimum width=170mm, minimum height=14mm, inner sep=6pt,
    font=\sffamily\small\color{gray80}, align=center]
    at ([yshift=-8.5cm]current page.center) {%
    {\color{success}\faLightbulb}\;
    忽略空气密度（$\rho_a\approx 0$）时简化为：$\frac{dh}{dt} = \frac{A_1 V_1 - A_2 V_2}{A_0}$
    \quad 进多出少则液面升高
  };
\end{tikzpicture}
\pagemark{6/8}

\clearpage
% ============================================================
%  PAGE 7: 避坑指南 (DOs vs DON'Ts)
% ============================================================
\begin{tikzpicture}[remember picture, overlay]
  \fill[bgPage] (current page.south west) rectangle (current page.north east);

  % 页眉
  \fill[danger] (current page.north west) rectangle ([yshift=-2cm]current page.north east);
  \node[anchor=west,font=\sffamily\LARGE\bfseries\color{white}]
    at ([xshift=2cm,yshift=-1cm]current page.north west)
    {\faExclamationTriangle\; 考前避坑指南：五大易错点};

  \tikzset{
    errbox/.style={
      rectangle, rounded corners=5pt, minimum width=80mm, minimum height=18mm,
      inner sep=6pt, align=left, font=\sffamily\small, text width=72mm
    }
  }

  % ---- 错误 1 ----
  \node[errbox, fill=danger!8, draw=danger!40, line width=0.6pt] (e1)
    at ([xshift=-4.5cm,yshift=4.2cm]current page.center) {%
    {\color{danger}\bfseries \faTimesCircle\; 系统 = 控制体}\\[2pt]
    {\color{gray80}把两者搞混，胡乱套用公式}
  };
  \node[errbox, fill=success!8, draw=success!40, line width=0.6pt] (c1)
    at ([xshift=4.5cm,yshift=4.2cm]current page.center) {%
    {\color{success}\bfseries \faCheckCircle\; 明确区分视角}\\[2pt]
    {\color{gray80}系统=跟着物质走；控制体=钉在空间}
  };
  \draw[-{Stealth[length=4pt]},gray40,line width=1pt] (e1) -- (c1);

  % ---- 错误 2 ----
  \node[errbox, fill=danger!8, draw=danger!40, line width=0.6pt] (e2)
    at ([xshift=-4.5cm,yshift=1.6cm]current page.center) {%
    {\color{danger}\bfseries \faTimesCircle\; 通量正负号搞反}\\[2pt]
    {\color{gray80}流入流出方向弄错，答案差个负号}
  };
  \node[errbox, fill=success!8, draw=success!40, line width=0.6pt] (c2)
    at ([xshift=4.5cm,yshift=1.6cm]current page.center) {%
    {\color{success}\bfseries \faCheckCircle\; 牢记外法向约定}\\[2pt]
    {\color{gray80}$\vec{n}$朝外：流出为正，流入为负}
  };
  \draw[-{Stealth[length=4pt]},gray40,line width=1pt] (e2) -- (c2);

  % ---- 错误 3 ----
  \node[errbox, fill=danger!8, draw=danger!40, line width=0.6pt] (e3)
    at ([xshift=-4.5cm,yshift=-1cm]current page.center) {%
    {\color{danger}\bfseries \faTimesCircle\; $\nabla\!\cdot\!\vec{v}\!=\!0$ 用于理想流体}\\[2pt]
    {\color{gray80}``理想流体''与``不可压缩''混淆}
  };
  \node[errbox, fill=success!8, draw=success!40, line width=0.6pt] (c3)
    at ([xshift=4.5cm,yshift=-1cm]current page.center) {%
    {\color{success}\bfseries \faCheckCircle\; $\nabla\!\cdot\!\vec{v}\!=\!0$ 是不可压条件}\\[2pt]
    {\color{gray80}理想流体也可以可压缩}
  };
  \draw[-{Stealth[length=4pt]},gray40,line width=1pt] (e3) -- (c3);

  % ---- 错误 4 ----
  \node[errbox, fill=danger!8, draw=danger!40, line width=0.6pt] (e4)
    at ([xshift=-4.5cm,yshift=-3.6cm]current page.center) {%
    {\color{danger}\bfseries \faTimesCircle\; 非均匀流速直接用 $AV$}\\[2pt]
    {\color{gray80}抛物线分布还用平均速度$\times$面积}
  };
  \node[errbox, fill=success!8, draw=success!40, line width=0.6pt] (c4)
    at ([xshift=4.5cm,yshift=-3.6cm]current page.center) {%
    {\color{success}\bfseries \faCheckCircle\; 用积分 $Q\!=\!\int v\,dA$}\\[2pt]
    {\color{gray80}层流抛物线分布：$Q=\frac{\pi R^2 v_0}{2}$}
  };
  \draw[-{Stealth[length=4pt]},gray40,line width=1pt] (e4) -- (c4);

  % ---- 错误 5 ----
  \node[errbox, fill=danger!8, draw=danger!40, line width=0.6pt] (e5)
    at ([xshift=-4.5cm,yshift=-6.2cm]current page.center) {%
    {\color{danger}\bfseries \faTimesCircle\; 水箱题忘了空气}\\[2pt]
    {\color{gray80}液面上有空气时只考虑了水}
  };
  \node[errbox, fill=success!8, draw=success!40, line width=0.6pt] (c5)
    at ([xshift=4.5cm,yshift=-6.2cm]current page.center) {%
    {\color{success}\bfseries \faCheckCircle\; 水和空气一起算}\\[2pt]
    {\color{gray80}封闭水箱出现 $(\rho-\rho_a)$ 项}
  };
  \draw[-{Stealth[length=4pt]},gray40,line width=1pt] (e5) -- (c5);

\end{tikzpicture}
\pagemark{7/8}

\clearpage
% ============================================================
%  PAGE 8: 知识树速查
% ============================================================
\begin{tikzpicture}[remember picture, overlay]
  \fill[bgPage] (current page.south west) rectangle (current page.north east);

  % 页眉
  \fill[purple] (current page.north west) rectangle ([yshift=-2cm]current page.north east);
  \node[anchor=west,font=\sffamily\LARGE\bfseries\color{white}]
    at ([xshift=2cm,yshift=-1cm]current page.north west)
    {\faSitemap\; 全章知识树速查};

  % 根节点
  \begin{scope}[shift={([xshift=-6cm,yshift=3.5cm]current page.center)}]
    \node[rectangle, rounded corners=6pt, fill=purple, text=white, minimum height=14mm,
      font=\sffamily\bfseries\large, inner xsep=12pt,
      blur shadow={shadow blur steps=4, shadow opacity=12}] (root) at (0,0) {第四章};

    % L1 nodes
    \node[rectangle, rounded corners=4pt, fill=primary, text=white,
      minimum height=10mm, font=\sffamily\bfseries, inner xsep=8pt] (l1a) at (5,4) {基本概念};
    \node[rectangle, rounded corners=4pt, fill=ocean, text=white,
      minimum height=10mm, font=\sffamily\bfseries, inner xsep=8pt] (l1b) at (5,1) {雷诺输运方程};
    \node[rectangle, rounded corners=4pt, fill=success, text=white,
      minimum height=10mm, font=\sffamily\bfseries, inner xsep=8pt] (l1c) at (5,-2) {连续性方程};
    \node[rectangle, rounded corners=4pt, fill=accent, text=white,
      minimum height=10mm, font=\sffamily\bfseries, inner xsep=8pt] (l1d) at (5,-5) {工程应用};

    % Connections root -> L1
    \draw[purple!50,thick,-{Stealth},out=0,in=180] (root.east) to (l1a.west);
    \draw[purple!50,thick,-{Stealth},out=0,in=180] (root.east) to (l1b.west);
    \draw[purple!50,thick,-{Stealth},out=0,in=180] (root.east) to (l1c.west);
    \draw[purple!50,thick,-{Stealth},out=0,in=180] (root.east) to (l1d.west);

    % L2 nodes for 基本概念
    \node[rectangle, rounded corners=3pt, fill=primaryLight!20, draw=primaryLight!50,
      font=\sffamily\small, inner sep=4pt] (l2a1) at (11,5.5) {系统 vs 控制体};
    \node[rectangle, rounded corners=3pt, fill=primaryLight!20, draw=primaryLight!50,
      font=\sffamily\small, inner sep=4pt] (l2a2) at (11,4) {控制面分类};
    \node[rectangle, rounded corners=3pt, fill=primaryLight!20, draw=primaryLight!50,
      font=\sffamily\small, inner sep=4pt] (l2a3) at (11,2.5) {广延量 vs 强度量};
    \draw[primaryLight,thick,-{Stealth},out=0,in=180] (l1a.east) to (l2a1.west);
    \draw[primaryLight,thick,-{Stealth},out=0,in=180] (l1a.east) to (l2a2.west);
    \draw[primaryLight,thick,-{Stealth},out=0,in=180] (l1a.east) to (l2a3.west);

    % L2 for RTT
    \node[rectangle, rounded corners=3pt, fill=oceanLight!20, draw=oceanLight!50,
      font=\sffamily\small, inner sep=4pt] (l2b1) at (11,1) {公式与物理含义};
    \draw[oceanLight,thick,-{Stealth},out=0,in=180] (l1b.east) to (l2b1.west);

    % L2 for 连续性
    \node[rectangle, rounded corners=3pt, fill=successLight!20, draw=successLight!50,
      font=\sffamily\small, inner sep=4pt] (l2c1) at (11,-1) {积分形式};
    \node[rectangle, rounded corners=3pt, fill=successLight!20, draw=successLight!50,
      font=\sffamily\small, inner sep=4pt] (l2c2) at (11,-2.2) {微分形式};
    \node[rectangle, rounded corners=3pt, fill=successLight!20, draw=successLight!50,
      font=\sffamily\small, inner sep=4pt] (l2c3) at (11,-3.4) {$\nabla\cdot\vec{v}=0$ 简化};
    \draw[successLight,thick,-{Stealth},out=0,in=180] (l1c.east) to (l2c1.west);
    \draw[successLight,thick,-{Stealth},out=0,in=180] (l1c.east) to (l2c2.west);
    \draw[successLight,thick,-{Stealth},out=0,in=180] (l1c.east) to (l2c3.west);

    % L2 for 工程应用
    \node[rectangle, rounded corners=3pt, fill=accentLight!20, draw=accentLight!50,
      font=\sffamily\small, inner sep=4pt] (l2d1) at (11,-4.5) {管道流 $A_1V_1=A_2V_2$};
    \node[rectangle, rounded corners=3pt, fill=accentLight!20, draw=accentLight!50,
      font=\sffamily\small, inner sep=4pt] (l2d2) at (11,-5.7) {速度分布求流量};
    \node[rectangle, rounded corners=3pt, fill=accentLight!20, draw=accentLight!50,
      font=\sffamily\small, inner sep=4pt] (l2d3) at (11,-6.9) {水箱液面变化率};
    \draw[accentLight,thick,-{Stealth},out=0,in=180] (l1d.east) to (l2d1.west);
    \draw[accentLight,thick,-{Stealth},out=0,in=180] (l1d.east) to (l2d2.west);
    \draw[accentLight,thick,-{Stealth},out=0,in=180] (l1d.east) to (l2d3.west);
  \end{scope}

  % 底部总结
  \node[rectangle, rounded corners=6pt, fill=purple!8, draw=purple!30,
    line width=0.8pt, minimum width=170mm, minimum height=20mm,
    inner sep=8pt, align=center, font=\sffamily\small\color{gray80}]
    at ([yshift=-9cm]current page.center) {%
    {\color{purple}\bfseries \faStar\; 一句话总结}\quad
    雷诺输运方程是``翻译工具''，连续性方程是``质量守恒''，工程上最常用 $A_1V_1=A_2V_2$
  };
\end{tikzpicture}
\pagemark{8/8}

\end{document}
